\chapter{Complexity — Sequential Circuit}
%%% generated by the blueprint pipeline from TCSlib.Complexity.CircuitComplexity.SequentialCircuit
\section{Overview}

Declarations of the Lean module \texttt{TCSlib.Complexity.CircuitComplexity.SequentialCircuit}.

\section{Declarations}

\begin{definition}[Sequential Boolean circuit]
\lean{BoolCircuit.SeqCircuit}
\label{BoolCircuit.SeqCircuit}
\leanok
\uses{BoolCircuit.DAGCircuit}
A sequential Boolean circuit with $n$ input bits and $r$ registers consists of a
combinational directed acyclic circuit with $n+r$ sources, a choice of one vertex of
that circuit to supply the next value of each register, and an initial Boolean value for
each register. Each selected vertex must be among the circuit's $n+r$ sources or its
gates.
\end{definition}

\begin{definition}[Sequential circuit size]
\lean{BoolCircuit.SeqCircuit.size}
\label{BoolCircuit.SeqCircuit.size}
\leanok
For a sequential Boolean circuit with $r$ registers and a combinational circuit
containing $g$ gates, define its size to be $r+g$.
\end{definition}

\begin{definition}[Sequential circuit state]
\lean{BoolCircuit.SeqCircuit.state}
\label{BoolCircuit.SeqCircuit.state}
\leanok
\uses{BoolCircuit.DAGCircuit.values}
For a sequential Boolean circuit $M$ with $n$ inputs and $r$ registers, and an input
assignment $x\colon\{0,\ldots,n-1\}\to\{\mathrm{false},\mathrm{true}\}$, define its
state $s_t\colon\{0,\ldots,r-1\}\to\{\mathrm{false},\mathrm{true}\}$ at each time
$t\in\mathbb{N}$ recursively: $s_0$ is the initial register state, and $s_{t+1}(j)$ is
the value at the next-value index specified for register $j$ when the combinational
circuit is evaluated on the input assignment followed by $s_t$, with any index outside
the resulting value list assigned value false.
\end{definition}

\begin{definition}[Sequential circuit output]
\lean{BoolCircuit.SeqCircuit.run}
\label{BoolCircuit.SeqCircuit.run}
\leanok
\uses{BoolCircuit.DAGCircuit.eval, BoolCircuit.SeqCircuit.state}
For a sequential Boolean circuit, a time $T\in\mathbb{N}$, and a Boolean assignment $x$
to its inputs, define the output at time $T$ by evaluating its combinational circuit on
$x$ together with the register state at time $T$.
\end{definition}

\begin{definition}[Vertex placement across time steps]
\lean{BoolCircuit.SeqCircuit.place}
\label{BoolCircuit.SeqCircuit.place}
\leanok
\uses{BoolCircuit.SeqCircuit.size}
For a sequential Boolean circuit of size $s$, define the placement of a vertex index $a$
at time step $k$ to be $a$ if $a<n$, and $a+ks$ otherwise, where $n$ is the number of
inputs.
\end{definition}

\begin{definition}[Gate in the unrolled circuit]
\lean{BoolCircuit.SeqCircuit.unrollGate}
\label{BoolCircuit.SeqCircuit.unrollGate}
\leanok
\uses{BoolCircuit.DAGGate, BoolCircuit.SeqCircuit.place, BoolCircuit.SeqCircuit.size}
Let $M$ be a sequential Boolean circuit with $n$ inputs and $r$ registers, and let $p$
be a nonnegative integer. Write $s$ for the circuit size, $q=\lfloor p/s\rfloor$ for the
block number, and $j=p\bmod s$ for the position within the block. If $j<r$ and $q=0$,
the gate is an argument-free AND gate when register $j$ is initially true and an
argument-free OR gate otherwise; if $j<r$ and $q>0$, it is an AND gate whose sole
argument is the vertex for register $j$'s next value placed in block $q-1$. If $j\ge r$,
take the combinational gate at index $j-r$ (using an argument-free AND gate if that
index is out of range); the resulting gate has that gate's kind and its arguments placed
in block $q$.
\end{definition}

\begin{theorem}[Bound on placed vertices]
\lean{BoolCircuit.SeqCircuit.place_lt}
\label{BoolCircuit.SeqCircuit.place_lt}
\leanok
\uses{BoolCircuit.SeqCircuit.place, BoolCircuit.SeqCircuit.size}
\difficulty{3}
Let a sequential circuit have $n$ inputs, and let its combinational circuit have $r$
register sources and $g$ gates. For every natural number $k$, if $a<n+r+g$, then the
vertex obtained by placing $a$ at stage $k$ is less than $n+(k+1)(r+g)$.
\end{theorem}

\begin{theorem}[Argument bound in an unrolled circuit]
\lean{BoolCircuit.SeqCircuit.args_unrollGate_lt}
\label{BoolCircuit.SeqCircuit.args_unrollGate_lt}
\leanok
\uses{BoolCircuit.SeqCircuit.place, BoolCircuit.SeqCircuit.place_lt, BoolCircuit.SeqCircuit.size, BoolCircuit.SeqCircuit.unrollGate}
\difficulty{5}
Let $M$ be a sequential circuit with $n$ inputs and $r$ registers, and let $p$ be a
natural number. Every vertex index appearing among the arguments of the gate at position
$p$ in the unrolled circuit is less than $n+p$.
\end{theorem}

\begin{definition}[Finite-time circuit unrolling]
\lean{BoolCircuit.SeqCircuit.unroll}
\label{BoolCircuit.SeqCircuit.unroll}
\leanok
\uses{BoolCircuit.DAGCircuit, BoolCircuit.SeqCircuit.args_unrollGate_lt, BoolCircuit.SeqCircuit.place, BoolCircuit.SeqCircuit.place_lt, BoolCircuit.SeqCircuit.size, BoolCircuit.SeqCircuit.unrollGate}
For a sequential Boolean circuit with $n$ inputs and any $T\in\mathbb{N}$, define a
directed acyclic circuit with $n$ sources by listing, in order, the gates at positions
$0$ through $(T+1)s-1$ from the circuit's unrolling, where $s$ is the sequential
circuit's size. Its output vertex is the placement at time step $T$ of the original
combinational circuit's output vertex.
\end{definition}

\begin{theorem}[Quotient and remainder within a block]
\lean{BoolCircuit.SeqCircuit.div_mod_block}
\label{BoolCircuit.SeqCircuit.div_mod_block}
\leanok
\difficulty{3}
For natural numbers $S,k,m$ with $m<S$, the quotient and remainder on dividing $Sk+m$ by
$S$ are respectively $k$ and $m$.
\end{theorem}

\begin{theorem}[Gate at a position in the unrolled circuit]
\lean{BoolCircuit.SeqCircuit.getElem_unroll_gates}
\label{BoolCircuit.SeqCircuit.getElem_unroll_gates}
\leanok
\uses{BoolCircuit.SeqCircuit.unroll, BoolCircuit.SeqCircuit.unrollGate}
\difficulty{1}
Let a sequential circuit be unrolled for $T$ steps, and let $p$ be an index smaller than
the number of gates in the resulting DAG circuit. The gate at position $p$ is the gate
specified by the unrolling rule at position $p$.
\end{theorem}

\begin{theorem}[Number of gates in the unrolled circuit]
\lean{BoolCircuit.SeqCircuit.length_unroll_gates}
\label{BoolCircuit.SeqCircuit.length_unroll_gates}
\leanok
\uses{BoolCircuit.SeqCircuit.size, BoolCircuit.SeqCircuit.unroll}
\difficulty{1}
Let $M$ be a sequential circuit with size $s$, and let $T$ be a natural number. The
circuit obtained by unrolling $M$ for $T$ steps has exactly $(T+1)s$ gates.
\end{theorem}

\begin{theorem}[Values within an unrolled stage]
\lean{BoolCircuit.SeqCircuit.unroll_block}
\label{BoolCircuit.SeqCircuit.unroll_block}
\leanok
\uses{BoolCircuit.DAGCircuit.values, BoolCircuit.DAGCircuit.values_getD_gate, BoolCircuit.DAGCircuit.values_getD_input, BoolCircuit.DAGGate.eval_remap, BoolCircuit.SeqCircuit.div_mod_block, BoolCircuit.SeqCircuit.getElem_unroll_gates, BoolCircuit.SeqCircuit.length_unroll_gates, BoolCircuit.SeqCircuit.place, BoolCircuit.SeqCircuit.size, BoolCircuit.SeqCircuit.state, BoolCircuit.SeqCircuit.unroll, BoolCircuit.SeqCircuit.unrollGate}
\difficulty{6}
Let a sequential circuit be unrolled for $T$ steps on input $x$, and let $k\le T$.
Suppose that at stage $k$ every register vertex has the value prescribed by the
sequential circuit's state at time $k$. Then for every vertex index $a<n+r+g$, the value
at its placed vertex in the unrolled circuit equals the value at index $a$ when the
combinational circuit is evaluated on the input together with the time-$k$ register
state; here $g$ is the number of combinational gates.
\end{theorem}

\begin{theorem}[Register values in the unrolled circuit]
\lean{BoolCircuit.SeqCircuit.unroll_reg}
\label{BoolCircuit.SeqCircuit.unroll_reg}
\leanok
\uses{BoolCircuit.DAGCircuit.values, BoolCircuit.DAGCircuit.values_getD_gate, BoolCircuit.DAGGate, BoolCircuit.DAGGate.eval, BoolCircuit.SeqCircuit.div_mod_block, BoolCircuit.SeqCircuit.getElem_unroll_gates, BoolCircuit.SeqCircuit.length_unroll_gates, BoolCircuit.SeqCircuit.place, BoolCircuit.SeqCircuit.size, BoolCircuit.SeqCircuit.unroll, BoolCircuit.SeqCircuit.unrollGate}
\difficulty{5}
Let $M$ be a sequential Boolean circuit with $n$ inputs and $r$ registers, let $T$ be a
nonnegative integer, and fix an input assignment
$x\colon\{0,\ldots,n-1\}\to\{\mathrm{false},\mathrm{true}\}$. For each time $k$ with
$k\le T$ and each register $j$, the value at the unrolled circuit vertex assigned to
register $j$ at time $k$ equals the evaluation, on the values of the unrolled circuit,
of the gate with no inputs whose type is AND or OR according to whether register $j$ is
initially true or false when $k=0$, and of the one-input AND gate reading register $j$'s
next-state source at time $k-1$ when $k>0$.
\end{theorem}

\begin{theorem}[Values in the unrolled circuit]
\lean{BoolCircuit.SeqCircuit.unroll_values}
\label{BoolCircuit.SeqCircuit.unroll_values}
\leanok
\uses{BoolCircuit.DAGCircuit.values, BoolCircuit.DAGGate.eval, BoolCircuit.SeqCircuit.place, BoolCircuit.SeqCircuit.state, BoolCircuit.SeqCircuit.unroll, BoolCircuit.SeqCircuit.unroll_block, BoolCircuit.SeqCircuit.unroll_reg}
\difficulty{4}
Let $M$ be a sequential Boolean circuit with $n$ inputs and $r$ registers, let $T$ be a
nonnegative integer, and fix an input assignment
$x\colon\{0,\ldots,n-1\}\to\{\mathrm{false},\mathrm{true}\}$. For every $k\le T$ and
every vertex index $a<n+r+g$, where $g$ is the number of gates in the combinational
part, the value at the vertex assigned by the unrolling to index $a$ at time $k$ equals
the value at index $a$ of the combinational circuit evaluated on the concatenation of
$x$ and the state of $M$ at time $k$.
\end{theorem}

\begin{theorem}[Evaluation of the unrolled circuit]
\lean{BoolCircuit.SeqCircuit.unroll_eval}
\label{BoolCircuit.SeqCircuit.unroll_eval}
\leanok
\uses{BoolCircuit.DAGCircuit.eval, BoolCircuit.SeqCircuit.run, BoolCircuit.SeqCircuit.unroll, BoolCircuit.SeqCircuit.unroll_values}
\difficulty{1}
For every sequential Boolean circuit, every natural number $T$, and every input $x$, the
output of the DAG circuit obtained by unrolling for $T$ steps equals the sequential
circuit's output after $T$ steps.
\end{theorem}

\begin{theorem}[Size of the unrolled circuit]
\lean{BoolCircuit.SeqCircuit.unroll_size}
\label{BoolCircuit.SeqCircuit.unroll_size}
\leanok
\uses{BoolCircuit.DAGCircuit.size, BoolCircuit.SeqCircuit.length_unroll_gates, BoolCircuit.SeqCircuit.size, BoolCircuit.SeqCircuit.unroll}
\difficulty{1}
Let $M$ be a sequential Boolean circuit with $n$ inputs, and let its size be the number
of registers plus the number of gates in its combinational part. For every nonnegative
integer $T$, the size of its unrolling through time $T$ is $n+(T+1)$ times the size of
$M$.
\end{theorem}

\begin{theorem}[Injectivity of vertex placement]
\lean{BoolCircuit.SeqCircuit.place_injective}
\label{BoolCircuit.SeqCircuit.place_injective}
\leanok
\uses{BoolCircuit.SeqCircuit.place}
\difficulty{4}
For a sequential circuit and any natural number $k$, the vertex placement map for stage
$k$ is injective.
\end{theorem}

\begin{theorem}[Cases for an unrolled gate]
\lean{BoolCircuit.SeqCircuit.unrollGate_cases}
\label{BoolCircuit.SeqCircuit.unrollGate_cases}
\leanok
\uses{BoolCircuit.SeqCircuit.place, BoolCircuit.SeqCircuit.size, BoolCircuit.SeqCircuit.unrollGate}
\difficulty{5}
For any sequential circuit and any natural number $p$, either the gate specified at
position $p$ has at most one argument and is not a NOT gate, or it equals a gate from
the combinational circuit with its arguments placed at stage $\lfloor p/s\rfloor$, where
$s$ is the sequential circuit's size.
\end{theorem}

\begin{theorem}[Well-formedness of the unrolled circuit]
\lean{BoolCircuit.SeqCircuit.unroll_isWellFormed}
\label{BoolCircuit.SeqCircuit.unroll_isWellFormed}
\leanok
\uses{BoolCircuit.DAGCircuit.IsWellFormed, BoolCircuit.SeqCircuit.place_injective, BoolCircuit.SeqCircuit.unroll, BoolCircuit.SeqCircuit.unrollGate_cases}
\difficulty{4}
Let $M$ be a sequential Boolean circuit with combinational part $C$, and let $T$ be a
nonnegative integer. If each gate of $C$ has distinct input vertices and every NOT gate
of $C$ has exactly one input, then each gate in the directed acyclic circuit obtained by
unrolling $M$ through time $T$ has distinct input vertices, and every NOT gate in that
circuit has exactly one input.
\end{theorem}

\begin{theorem}[Fan-in of the unrolled circuit]
\lean{BoolCircuit.SeqCircuit.unroll_isFaninTwo}
\label{BoolCircuit.SeqCircuit.unroll_isFaninTwo}
\leanok
\uses{BoolCircuit.DAGCircuit.IsFaninTwo, BoolCircuit.SeqCircuit.unroll, BoolCircuit.SeqCircuit.unrollGate_cases, BoolCircuit.SeqCircuit.unroll_isWellFormed}
\difficulty{4}
Let $M$ be a sequential Boolean circuit with combinational part $C$, and let $T$ be a
nonnegative integer. If every gate of $C$ has distinct input vertices, every NOT gate of
$C$ has exactly one input, and every gate of $C$ has at most two inputs, then the
directed acyclic circuit obtained by unrolling $M$ through time $T$ has distinct input
vertices at each gate, exactly one input at every NOT gate, and at most two inputs at
every gate.
\end{theorem}

\begin{theorem}[DAG simulation of a sequential circuit]
\lean{BoolCircuit.SeqCircuit.exists_dagCircuit_of_run}
\label{BoolCircuit.SeqCircuit.exists_dagCircuit_of_run}
\leanok
\uses{BoolCircuit.DAGCircuit, BoolCircuit.DAGCircuit.IsFaninTwo, BoolCircuit.DAGCircuit.IsWellFormed, BoolCircuit.DAGCircuit.eval, BoolCircuit.DAGCircuit.size, BoolCircuit.SeqCircuit, BoolCircuit.SeqCircuit.run, BoolCircuit.SeqCircuit.size, BoolCircuit.SeqCircuit.unroll, BoolCircuit.SeqCircuit.unroll_eval, BoolCircuit.SeqCircuit.unroll_isFaninTwo, BoolCircuit.SeqCircuit.unroll_isWellFormed, BoolCircuit.SeqCircuit.unroll_size}
\difficulty{4}
Let a Boolean sequential circuit have $n$ input bits, $r$ registers, a combinational DAG
on its inputs and register contents, a next value for each register, and an initial
register state. For every integer $T\ge 1$, there is a DAG circuit on the same $n$
inputs whose output agrees with the sequential circuit's output after $T$ steps, and
whose size is at most $2((r+g)T+n)$, where $g$ is the number of gates in the
combinational DAG. If the combinational DAG is well formed, so is the resulting DAG; if
it is well formed and every gate has at most two arguments, then the resulting DAG has
those same properties.
\end{theorem}
